\section{Phase 2 Post-Deployment Interview Protocols}

Phase 2 used separate semi-structured interview guides for older adults and family caregivers after the two-week prototype deployment. Interviewers first asked participants to describe their experiences in their own terms and then used probes to reconstruct specific episodes. Questions about features participants had used were distinguished from questions about design mechanisms that had not been implemented. Unimplemented mechanisms were introduced hypothetically and were treated as design judgments rather than accounts of use.

\subsection{Appendix B1: Older-Adult Post-Deployment Interview Protocol}

\subsubsection{Section 1: Overall Experience and Everyday Use}

\begin{enumerate}
  \item \textbf{Can you tell me how using the system went for you over the past two weeks?}

  \textit{Probe: What became part of your usual routine?} \textit{Probe: Was there anything you particularly liked, disliked, or stopped paying attention to?}
  \item \textbf{Can you walk me through a recent day when you used the system, from the first reminder onward?}

  \textit{Probe: Where were you? What did the system do? What did you do next?}
  \item \textbf{Did using the system change anything about how you managed your medicines?}

  \textit{Probe: What, if anything, did you do differently from before?}
  \item \textbf{Were there times when the system did not fit what was happening that day?}

  \textit{Probe: Tell me about the most recent time.}
\end{enumerate}

\subsubsection{Section 2: Reminders, Confirmation, and Family Involvement}

\begin{enumerate}[resume]
  \item \textbf{What usually happened when a medicine reminder appeared?}

  \textit{Probe: What did you normally do after hearing or seeing it?}
  \item \textbf{Were there occasions when you did not confirm a dose?}

  \textit{Probe: Tell me about a specific occasion. Why was it left unconfirmed?}
  \item \textbf{What happened when a dose remained unconfirmed?}

  \textit{Probe: What did the system do next?}
  \item \textbf{Did the system ever ask whether a family member should be involved?}

  \textit{Probe: Tell me about the most recent time you remember.}
  \item \textbf{How did you decide whether to agree or decline when it asked?}

  \textit{Probe: What was happening in that situation?}
  \item \textbf{Were there situations when you preferred to handle the matter yourself?}

  \textit{Probe: What made those situations different?}
  \item \textbf{Were there situations when involving a family member was useful to you?}

  \textit{Probe: What made their involvement useful in that case?}
  \item \textbf{After a family member became involved, what happened next?}

  \textit{Probe: Did they call, ask you something, come to you, or respond in another way?}
\end{enumerate}

\subsubsection{Section 3: Who the System Was Serving}

\begin{enumerate}[resume]
  \item \textbf{Did the system ever say who it was serving or working with?}

  \textit{Probe: What do you remember it saying?}
  \item \textbf{What did that announcement mean to you, if anything?}

  \textit{Probe: Did you usually notice it, or did it become part of the background?}
  \item \textbf{Did there seem to be times when the system was mainly helping you and other times when it was bringing your family into the situation?}

  \textit{Probe: Can you describe one of each?}
  \item \textbf{Was there ever a time when you wanted the system to involve someone differently from what it did?}

  \textit{Probe: What would you have preferred?}
  \item \textbf{Did you ever want to stop, delay, or change what the system was about to do?}

  \textit{Probe: What happened?}
\end{enumerate}

\subsubsection{Section 4: Family Checking and Medication Records}

\begin{enumerate}[resume]
  \item \textbf{Before using the system, how did your family usually know whether you had taken your medicines?}
  \item \textbf{During the two weeks, did anything change about how family members checked or asked about your medicines?}

  \textit{Probe: More often, less often, or in a different way?}
  \item \textbf{How did you feel about the way family members used the information from the system?}

  \textit{Probe: Can you give me a specific example?}
  \item \textbf{Did the medication record, logbook, streak, or heat map ever come up in a conversation with someone in your family?}

  \textit{Probe: Who brought it up, and what was said?}
  \item \textbf{Did you pay attention to the streak or medication record yourself?}

  \textit{Probe: Did it affect anything you did?}
  \item \textbf{Who do you think should be able to see medication records like these?}

  \textit{Probe: Should that remain the same all the time, or should it be possible to change?}
  \item \textbf{Was there ever a difference between what the system recorded and what you or someone in your family believed had happened?}

  \textit{Probe: How was that handled?}
\end{enumerate}

\subsubsection{Section 5: Trust, Reliance, and Breakdown}

\begin{enumerate}[resume]
  \item \textbf{Think back to the first few days of using the system. How did you decide whether you could rely on it?}

  \textit{Probe: Did you continue checking anything yourself?}
  \item \textbf{Did your way of using the system change as the two weeks went on?}

  \textit{Probe: What led to that change?}
  \item \textbf{Did the system ever get something wrong or behave differently from what you expected?}

  \textit{Probe: What did you do afterward?}
  \item \textbf{If the system stopped working tomorrow, what would you do about your medicines?}

  \textit{Probe: Would anything be difficult to return to?}
  \item \textbf{Do you think using the system changed how much you relied on your own memory, family members, or other reminders?}

  \textit{Probe: In what way?}
\end{enumerate}

\subsubsection{Section 6: Proposed Design Mechanisms}

\textit{The mechanisms in this section were described as possibilities rather than features participants had used.}

\begin{enumerate}[resume]
  \item \textbf{Suppose you told the system not to involve a family member, but it believed the situation was becoming more concerning. What should it do?}

  \textit{Probe: Should your refusal always remain final, or are there circumstances when it should ask again?}
  \item \textbf{If it asked again, how should it explain why it was doing so?}
  \item \textbf{Suppose you did not respond to a reminder at all. What should the system assume?}

  \textit{Probe: Are there ordinary reasons you might not respond?}
  \item \textbf{How could the system distinguish an ordinary non-response from a situation that may need attention?}
  \item \textbf{Imagine that when you first started using the system, it stayed in a more limited mode until you became comfortable with it. Would that be useful?}

  \textit{Probe: What would you want it to do during that period?} \textit{Probe: Who should decide when that period ends?}
  \item \textbf{Imagine that family members could also have shared scores or progress information connected to medication support. What would you think about that?}

  \textit{Probe: What, if anything, should be shared?} \textit{Probe: Who should decide whether it is visible?}
\end{enumerate}

\subsubsection{Section 7: Closing Reflection}

\begin{enumerate}[resume]
  \item \textbf{Taking everything together, who did the system seem to be working for during these two weeks?}

  \textit{Probe: Did that seem the same in every situation?}
  \item \textbf{Can you describe a situation in which you and a family member might want different things from the system?}

  \textit{Probe: What should the system do in that situation?}
  \item \textbf{When, if ever, should the system move from helping you privately to involving someone else?}

  \textit{Probe: Who should decide that?}
  \item \textbf{Can you take me through one recent interaction with the system from beginning to end?}

  \textit{Probe: What happened before it started, what did the system do first, what happened next, who became involved, and how did it end?}
  \item \textbf{Would you want to continue using a system like this?}

  \textit{Probe: Under what conditions?}
  \item \textbf{If you could change one thing about the system, what would you change?}
  \item \textbf{Is there anything about your experience that we have not discussed but that you think is important?}
\end{enumerate}

\subsection{Appendix B2: Family-Caregiver Post-Deployment Interview Protocol}

\subsubsection{Section 1: Overall Experience and Household Use}

\begin{enumerate}
  \item \textbf{Can you tell me how the past two weeks with the system went from your perspective?}

  \textit{Probe: What, if anything, changed in your usual medication-related work?}
  \item \textbf{Can you walk me through a recent occasion when the system became relevant to you?}

  \textit{Probe: What happened first, and how did you become involved?}
  \item \textbf{Was there anything about the system that was useful for you?}

  \textit{Probe: Can you describe a specific occasion?}
  \item \textbf{Was there anything that did not fit how medication care normally works in your household?}

  \textit{Probe: Tell me about the most recent example.}
\end{enumerate}

\subsubsection{Section 2: Reminders, Family Involvement, and Permission}

\begin{enumerate}[resume]
  \item \textbf{When a dose was due, what did you normally know about it?}

  \textit{Probe: Did the system contact you directly, or did you usually learn about it another way?}
  \item \textbf{Did the system ever involve you after a dose remained unconfirmed?}

  \textit{Probe: Tell me about a specific occasion.}
  \item \textbf{What happened once you were notified?}

  \textit{Probe: What did you do, and how did the older adult respond?}
  \item \textbf{Were there occasions when you expected to be involved but were not?}

  \textit{Probe: What made you expect involvement in that situation?}
  \item \textbf{Did you know whether the older adult could decline involving you?}

  \textit{Probe: What did you think about that arrangement?}
  \item \textbf{Were there situations when you thought the older adult should handle the matter without you?}

  \textit{Probe: What made those situations different?}
  \item \textbf{Were there situations when you thought family involvement was particularly important?}

  \textit{Probe: Why?}
\end{enumerate}

\subsubsection{Section 3: Who the System Was Serving}

\begin{enumerate}[resume]
  \item \textbf{Did the system ever say who it was serving or working with?}

  \textit{Probe: What do you remember it saying?}
  \item \textbf{What did you make of those announcements?}

  \textit{Probe: Did they affect how you understood your own role?}
  \item \textbf{Did the way the system involved you or the older adult seem to change across situations?}

  \textit{Probe: Tell me about a specific example.}
  \item \textbf{Was there a change you wanted in how the system involved family members that did not happen?}

  \textit{Probe: What would you have wanted instead?}
  \item \textbf{Was there anything the system told you that you were unsure the older adult knew had been shared?}

  \textit{Probe: What did you do in that situation?}
\end{enumerate}

\subsubsection{Section 4: Checking, Communication, and Medication Records}

\begin{enumerate}[resume]
  \item \textbf{Before the deployment, how did you usually find out whether medicines had been taken?}
  \item \textbf{During the deployment, did anything change about how you checked?}

  \textit{Probe: Did you ask more, less, or differently?}
  \item \textbf{Did anything change about your conversations with the older adult concerning medicines?}

  \textit{Probe: Was there a kind of conversation you had before that happened less often during deployment?}
  \item \textbf{Did the older adult ever say anything about how you were checking or responding to information from the system?}

  \textit{Probe: What do you remember them saying?}
  \item \textbf{Did you use the logbook, streak, or heat map to understand what had happened?}

  \textit{Probe: How did you use it, if at all?}
  \item \textbf{Did anyone in the family talk about the streak or medication record?}

  \textit{Probe: Who raised it, and what happened afterward?}
  \item \textbf{Did the system's record ever differ from what you or another family member thought had happened?}

  \textit{Probe: How was the difference handled?}
  \item \textbf{Who do you think should be able to see medication records like these?}

  \textit{Probe: Who should decide, and should that decision be changeable?}
\end{enumerate}

\subsubsection{Section 5: Trust, Reliance, and Breakdown}

\begin{enumerate}[resume]
  \item \textbf{Think back to the first few days. What did you do to decide whether you could rely on the system?}

  \textit{Probe: Did you continue checking anything independently?}
  \item \textbf{Did the way you relied on the system change during the two weeks?}

  \textit{Probe: What led to that change?}
  \item \textbf{Did the system ever get something wrong or behave differently from what you expected?}

  \textit{Probe: What did you do afterward?}
  \item \textbf{Do you still keep the older adult's medication schedule in your head in the same way as before?}

  \textit{Probe: Has anything changed?}
  \item \textbf{If the system stopped working tomorrow, what would happen in the household?}

  \textit{Probe: Walk me through how you would manage that day.}
\end{enumerate}

\subsubsection{Section 6: Proposed Design Mechanisms}

\textit{The mechanisms in this section were described as possibilities rather than features participants had used.}

\begin{enumerate}[resume]
  \item \textbf{Suppose the older adult declined family involvement, but the system believed the situation was becoming more concerning. What should happen next?}

  \textit{Probe: Should the system accept the decision, ask again, or do something else?}
  \item \textbf{Are there circumstances in which you think the system should continue asking even after someone has said no?}

  \textit{Probe: What limits should there be?}
  \item \textbf{Suppose nobody responded to a reminder. What should the system understand from that?}

  \textit{Probe: What ordinary situations could produce no response?}
  \item \textbf{How should it distinguish those situations from one that may need family attention?}
  \item \textbf{Imagine that during the first days of use, the system operated in a more limited mode while the household decided whether to trust it. Would that be useful?}

  \textit{Probe: What would you want to check during that period?} \textit{Probe: Who should decide when the system can do more?}
  \item \textbf{Imagine that family members could share scores or progress information connected to medication support. What would you think about that?}

  \textit{Probe: What should be visible, to whom, and who should decide?}
\end{enumerate}

\subsubsection{Section 7: Closing Reflection}

\begin{enumerate}[resume]
  \item \textbf{Taking everything together, who did the system seem to be working for during the deployment?}

  \textit{Probe: Did your answer change depending on the situation?}
  \item \textbf{Can you think of a situation in which what you wanted and what the older adult wanted were different?}

  \textit{Probe: What should the system do when that happens?}
  \item \textbf{When, if ever, should the system move from working directly with the older adult to involving the family?}

  \textit{Probe: Should it be able to make that change itself, or should someone authorize it?}
  \item \textbf{Can you take me through one recent interaction involving the system from beginning to end?}

  \textit{Probe: What happened beforehand, who was present, what did the system do, what did each person do, and how did it end?}
  \item \textbf{What do you think the older adult would say about that same interaction?}

  \textit{Record this separately from the caregiver's own account.}
  \item \textbf{Would you want the household to continue using a system like this?}

  \textit{Probe: Under what conditions?}
  \item \textbf{If you could change one thing about the system, what would you change?}
  \item \textbf{Is there anything about your experience or your medication-related work that we have not discussed but that you think is important?}
\end{enumerate}